\documentclass[12pt, a4paper]{article}

\usepackage[utf8]{inputenc}
\usepackage[T1]{fontenc}
\usepackage[english]{babel}
\usepackage{geometry}
\usepackage{listings}
\usepackage{xcolor}
\usepackage{fancyhdr}
\usepackage{parskip}

\geometry{margin=2.5cm}

\definecolor{codeBackground}{rgb}{0.95,0.95,0.95}
\definecolor{codeBorder}{rgb}{0.7,0.7,0.7}

\lstset{
    backgroundcolor=\color{codeBackground},
    basicstyle=\ttfamily\small,
    breaklines=true,
    frame=single,
    rulecolor=\color{codeBorder},
    xleftmargin=8pt,
    xrightmargin=8pt,
}

\pagestyle{fancy}
\fancyhf{}
\rhead{\textcolor{gray}{\small ClickWars -- multi-instance update}}
\lhead{\textcolor{gray}{\small Roulem Baitiche}}
\rfoot{\textcolor{gray}{\small Page \thepage}}
\lfoot{\textcolor{gray}{\small \today}}

\begin{document}

\begin{center}
    {\Large \textbf{ClickWars -- Multi-Instance Update}}\\[0.4em]
    {\small Roulem Baitiche \quad | \quad \today}
\end{center}

\vspace{1em}
\hrule
\vspace{1.2em}

\textbf{Task:} \textit{``1. Prepare the server to run safely in multiple parallel instances''}

\vspace{0.5em}
\hrule
\vspace{1em}

\section*{What was done}

The ClickWars server can now run as multiple instances in parallel behind a load balancer without game state going out of sync.

Previously, the game state (gauges, players, phase) lived only in the Node.js process memory. Two instances had no way to share it.

\textbf{Solution:} a new file \texttt{SharedStateStore.js} was added. It exposes two implementations of the same interface:

\begin{itemize}
    \item \texttt{MemoryStore} -- identical behaviour to the old code, no new dependencies. Used in development.
    \item \texttt{RedisStore} -- extends MemoryStore. Every state mutation is synced to other instances via Redis pub/sub. Each instance keeps a local cache for performance. A full state snapshot is saved to Redis every 2 seconds so a new instance can join an ongoing game.
\end{itemize}

\textbf{Other changes:}
\begin{itemize}
    \item \texttt{GameServer.js} -- uses \texttt{store.acquireLock()} for victory detection (prevents duplicate declarations) and for the bot loop (only one instance runs bots at a time).
    \item \texttt{websocket-server.js} -- each instance has a unique \texttt{INSTANCE\_ID}; client and bot IDs are prefixed to avoid collisions. The Redis pub/sub relay forwards WebSocket broadcasts to clients connected on other instances.
    \item No changes to the WebSocket protocol or the user interface.
\end{itemize}

\section*{Deployment prerequisite}

A Redis server must be reachable from each Node.js instance.

\begin{lstlisting}
# Install (Debian/Ubuntu)
apt install redis-server

# Verify
redis-cli ping   # -> PONG
\end{lstlisting}

\section*{Environment variables}

\begin{lstlisting}
REDIS_URL=redis://<host>:6379   # enables multi-instance mode
INSTANCE_ID=instance-1          # optional, auto-generated if absent
PORT=3000
\end{lstlisting}

Without \texttt{REDIS\_URL}, the server runs in memory mode (single instance, same behaviour as before).

\section*{Tests}

\begin{lstlisting}
# Unit tests (27 tests, no Redis required)
node tests.js

# Multi-instance tests (requires Redis)
node tests-multi-instance.js

# Visual demo (2 instances + Redis)
node demo-multi-instance.js
\end{lstlisting}

\section*{Nginx note}

For WebSocket connections, use \texttt{ip\_hash} in the upstream block so each client sticks to the same instance.

\begin{lstlisting}
upstream clickwars {
    ip_hash;
    server 127.0.0.1:3000;
    server 127.0.0.1:3001;
}
\end{lstlisting}

\end{document}
